\chapter{روال‌ها: شکستنِ برنامه به قطعه‌ها}%
\label{ch:functions}

هر سه آجرِ الگوریتم (توالی، انتخاب، تکرار) را داریم و از این پس هیچ
الگوریتمی از دستمان درنمی‌رود. اما با بزرگ‌تر شدنِ برنامه، مشکلِ تازه‌ای سر
برمی‌آورد: شلوغی. راهِ‌حلِ برنامه‌نویس‌ها از روزِ اول یکی بوده است: برنامهٔ
بزرگ را به قطعه‌های کوچکِ نام‌دار بشکن. به این قطعه‌ها \textbf{روال}
می‌گویند.

\section{روال چیست؟}

روال بسته‌ای از دستورهاست که نامی دارد و کارِ مشخصی می‌کند. یک بار
می‌نویسیدش و هر بار که لازم شد، فقط \emph{صدایش می‌زنید}. با یکی از
آشناترین‌ها شروع کنیم؛ همان «بزرگ‌ترین از میان دو عدد»:

\begin{salamfa}
روال بیشینه(الف : صحیح, ب : صحیح) : صحیح:
    اگر الف > ب:
        برگشت الف
    پایان
    برگشت ب
پایان

روال ریشه:
    سرچاپ بیشینه(7, 12)
    سرچاپ بیشینه(40, 3)
    سرچاپ بیشینه(5, 5)
پایان
\end{salamfa}

\begin{progout}
12
40
5
\end{progout}

ساختارِ تعریفِ روال:

\begin{center}
\code{روال} \;\textit{نام}\code{(}\textit{پارامترها}\code{)}\;:\;\textit{گونهِ خروجی}\;\code{:}
\quad\ldots\;بدنه\;\ldots\quad \code{پایان}
\end{center}

تکه‌تکه‌اش کنیم:

\begin{itemize}
  \item \code{بیشینه} نامِ روال است؛ همان قاعده‌های نام‌گذاریِ متغیرها
        اینجا هم برقرار است: نام باید بگوید روال \emph{چه می‌کند}.
  \item \code{(الف صحیح, ب صحیح)} دو \textbf{پارامتر} است: ورودی‌های
        روال، هر یک با نام و گونه.
  \item \code{صحیح}ِ پس از پرانتز، گونهِ \textbf{خروجی} است: این روال در
        پایان یک عددِ صحیح تحویل می‌دهد.
  \item \code{برگشت} خروجی را تحویل می‌دهد و کارِ روال را \emph{همان
        لحظه} تمام می‌کند. به همین دلیل در بدنهٔ بالا به \code{وگرنه}
        نیازی نبود: اگر \code{الف > ب} باشد، با \code{برگشت الف} کارِ
        روال تمام است و هرگز به خطِ \code{برگشت ب} نمی‌رسیم.
\end{itemize}

وقتی در روالِ ریشه می‌نویسیم \code{بیشینه(7, 12)}، می‌گوییم: «روالِ بیشینه
را با الفِ ۷ و بِ ۱۲ اجرا کن و هر چه برگرداند، همین‌جا بگذار.» به این کار
\textbf{واردسازی} می‌گویند. ترتیب مهم است: مقدارِ اول در پارامترِ اول
می‌نشیند و به همین ترتیب.

\section{چرا روال؟}

سه دلیلِ بزرگ:

\begin{itemize}
  \item \textbf{پرهیز از تکرار:} اگر جایی سه بار یک منطق را کپی کرده‌اید،
        وقتِ روال ساختن است. فردا که اشکالی پیدا شد، یک جا را درست
        می‌کنید نه سه جا را.
  \item \textbf{خوانایی:} \code{سرچاپ بیشینه(الف, ب)} را هر کسی می‌فهمد؛
        نسخهٔ باز شده‌اش را باید خط‌به‌خط خواند.
  \item \textbf{آزمون‌پذیری:} هر روال را می‌توان جدا آزمود. سه فراخوانیِ
        مثالِ بالا در واقع سه \emph{آزمایش} بودند؛ حالتِ سوم (دو مقدارِ
        برابر) همان حالتِ مرزی‌ای است که در تمرینِ فصلِ
        \ref{ch:algorithms} بررسی کردید.
\end{itemize}

\begin{tip}
قاعدهٔ طلایی: \emph{هر روال فقط یک کار بکند و آن را خوب بکند.} روالی که هم
حساب می‌کند، هم چاپ می‌کند و هم تصمیم می‌گیرد، نه خوانا است، نه
دوباره‌استفاده‌شدنی و نه آزمودنی. کوتاه نگهش دارید؛ اگر از حدودِ بیست خط
گذشت، احتمالاً دو کار در آن قاطی شده و باید دو روال شود.
\end{tip}

\section{روالِ بدونِ خروجی}

هر روالی چیزی برنمی‌گرداند؛ بعضی فقط \emph{کاری می‌کنند} (مثلاً چاپ). در
این حالت گونهِ خروجی را نمی‌نویسیم:

\begin{salamfa}
روال خط بکش:
    تکرار 30:
        چاپ "-"
    پایان
    سرچاپ ""
پایان

روال ریشه:
    خط بکش()
    سرچاپ "گزارش فروش"
    خط بکش()
پایان
\end{salamfa}

\begin{progoutfa}
------------------------------
گزارش فروش
------------------------------
\end{progoutfa}

دقت کنید که روالِ بی‌پارامتر را با پرانتزِ خالی صدا می‌زنیم:
\code{خط بکش()}.

\section{متغیرهای محلی}

متغیری که درونِ روال ساخته می‌شود، فقط درونِ همان روال وجود دارد و بیرون
از آن دیده نمی‌شود؛ به آن \textbf{متغیرِ محلی} می‌گویند. پارامترها هم
محلی‌اند. این جداسازی نعمت است: هر روال آشپزخانهٔ خودش را دارد و روال‌های
دیگر نمی‌توانند ندانسته ظرف‌هایش را جابه‌جا کنند. اگر روالی به داده‌ای
نیاز دارد، آن را از راهِ پارامتر بدهید و نتیجه را از راهِ \code{برگشت}
پس بگیرید؛ این تنها درِ ورودی و خروجیِ تمیزِ یک روال است.

\section{ترکیبِ روال‌ها}

خروجیِ یک روال می‌تواند ورودیِ روالی دیگر باشد؛ روال‌ها مثلِ قطعه‌های لگو در
هم چفت می‌شوند. بزرگ‌ترین از میانِ \emph{سه} عدد، بدونِ نوشتنِ حتی یک
مقایسهٔ تازه:

\begin{salamfa}
روال بیشینه(الف : صحیح, ب : صحیح) : صحیح:
    اگر الف > ب:
        برگشت الف
    پایان
    برگشت ب
پایان

روال بیشینه سه تایی(الف : صحیح, ب : صحیح, ج : صحیح) : صحیح:
    برگشت بیشینه(بیشینه(الف, ب), ج)
پایان

روال ریشه:
    سرچاپ بیشینه سه تایی(7, 12, 5)
پایان
\end{salamfa}

\begin{progout}
12
\end{progout}

خطِ \code{بیشینه(بیشینه(الف, ب), ج)} از درون به بیرون خوانده می‌شود: اول
بزرگ‌ترینِ الف و ب پیدا می‌شود، بعد حاصلش با ج سنجیده می‌شود. این «ساختن از
قطعه‌های کوچکِ مطمئن» جانِ کلامِ مهندسیِ نرم‌افزار است.

یک نمونهٔ کامل‌تر که شرط و حلقه و روال را با هم دارد؛ آزمونِ اول بودنِ یک
عدد (عددی که فقط بر ۱ و خودش بخش‌پذیر است):

\begin{salamfa}
روال اول است(عدد : صحیح) : منطقی:
    اگر عدد < 2:
        برگشت نادرست
    پایان
    ناپایا مقسوم علیه : صحیح = 2
    تا (مقسوم علیه * مقسوم علیه) <= عدد:
        اگر (عدد % مقسوم علیه) == 0:
            برگشت نادرست
        پایان
        مقسوم علیه = مقسوم علیه + 1
    پایان
    برگشت درست
پایان

روال ریشه:
    تکرار 2 تا 20 از n:
        اگر اول است(n):
            چاپ n
            چاپ " "
        پایان
    پایان
    سرچاپ ""
پایان
\end{salamfa}

\begin{progout}
2 3 5 7 11 13 17 19
\end{progout}

دو ریزه‌کاریِ این روال خواندنی است: نگهبانِ اولِ کار حالت‌های مرزی (اعدادِ
کمتر از ۲) را یک‌جا کنار می‌گذارد؛ و شرطِ
\code{(مقسوم علیه * مقسوم علیه) <= عدد} یعنی فقط تا جذرِ عدد می‌گردیم، چون
اگر عددی مقسوم‌علیهی بزرگ‌تر از جذرش داشته باشد، حتماً یکی کوچک‌تر از جذر
هم دارد. با همین یک شرط، آزمودنِ اعدادِ بزرگ به‌مراتب سریع‌تر می‌شود؛
نمونه‌ای کوچک از اینکه \emph{فکرِ الگوریتمی} چگونه به برنامهٔ بهتر تبدیل
می‌شود.

\section{تمرین‌ها}

\exercise روالِ \code{کمینه(الف صحیح, ب صحیح) صحیح} را بنویسید که
کوچک‌ترین را برگرداند، و با سه فراخوانی بیازمایید (فراموش نکنید: یکی از
آزمایش‌ها دو مقدارِ برابر باشد).

\exercise روالِ \code{قدر مطلق(عدد صحیح) صحیح} را بنویسید: اگر عدد منفی
بود قرینه‌اش را برگرداند، وگرنه خودش را.

\exercise روالِ \code{زوج است(عدد صحیح) منطقی} را بنویسید و سپس با کمکِ
آن (نه با نوشتنِ دوبارهٔ \code{\%})، اعدادِ زوجِ ۱ تا ۳۰ را چاپ کنید.

\exercise روالِ \code{توان(پایه صحیح, نما صحیح) صحیح} را با الگوی انباشتنِ
ضربی بنویسید؛ \code{توان(2, 10)} باید \code{1024} شود.

\exercise روالِ \code{شمارش مقسوم علیه ها(عدد صحیح) صحیح} را بنویسید که
تعدادِ مقسوم‌علیه‌های عدد را برگرداند. سپس در روالِ ریشه، با آن عددهای ۱
تا ۳۰ را بررسی کنید و هر عددی که دقیقاً دو مقسوم‌علیه دارد را چاپ کنید.
خروجی باید همان اعدادِ اول باشد. چرا؟
